% =====================================================================
\chapter{Conclusion}\label{ch:conclusion}
% =====================================================================

\section{Answers to the four questions}

\begin{problem}[Generality of low-level determinism]
The answer has three parts. The low level is nearly unique because it is
\emph{over-constrained}: the number of constraints is at least the number of design
degrees of freedom, and the solution set degenerates into a discrete point set
(Proposition \ref{prop:overconstrain}). The upper levels admit many solutions because
problem-level constraints are far fewer than the degrees of freedom, and the solution
set is a high-dimensional manifold. The upper levels nevertheless converge, because a
\emph{second set of constraints}, imposed by the human organisations that build and
maintain the system, compresses the manifold into a few sharp peaks
(Principle \ref{pr:org}).

The accurate answer is therefore this: convergence exists, but not towards the
engineering optimum; it is convergence towards a \emph{structure that a bounded
organisation can navigate}. What converges is carried by \emph{relations} rather than by
components (Principle \ref{pr:relations}); its strength decreases monotonically with
domain binding (Conjecture \ref{conj:strength}).
\end{problem}

\begin{problem}[Coercion by scale]
The conclusion is a \emph{conditional necessity} result. The independent variable that
decides whether a unified substrate is required is the product of the number of live
entities and the number of shared resources, not the line count
(Theorem \ref{th:integration}, Corollary \ref{cor:threshold}). The six functions
(identity and lifecycle, allocation and reclamation, scheduling, naming and addressing,
communication and synchronisation, and fault isolation and recovery) are necessary
conditions for the operational semantics to be well defined (Theorem \ref{th:six}). They
may be built or adopted from a host; the criterion is whether the invariant they carry
is domain-critical (Corollary \ref{cor:buildadopt}).
\end{problem}

\begin{problem}[Comprehension and design]
The conclusion is that comprehension is inverse compression and design is forward
compression, the two reading directions of one structure. The optimal descent depth of
comprehension is given by the stationarity condition of rate--distortion
(Proposition \ref{prop:stopping}); the measurable objective of design is the expected
read-set (Definition \ref{def:readset}), and minimising it requires the enforcement of
each invariant at a single site (Theorem \ref{th:centralise}). Whether an abstraction
leaks is measurable, not merely felt (Proposition \ref{prop:leakpair},
Corollary \ref{cor:diagnose}). For an AI the advantage lies not in reading more but in
probing faster (Proposition \ref{prop:active}).
\end{problem}

\begin{problem}[Correctness without review]
The conclusion is this: it is attainable that no human reviews the implementation line
by line, and unattainable that no human remains responsible for the expression of the
goal. ``Correctness'' must be decomposed into five dimensions
(Definition \ref{def:five-correctness}); the intermediate process must be defined as a
seven-stage verifiable refinement chain (Definition \ref{def:chain}); trust must be a
composite of the two modes of trust generation (Definition \ref{def:two-layer}); and
artefact binding is the fourth pillar and cannot be omitted (Eq. \eqref{eq:binding}). The
consolidated assurance condition (Theorem \ref{th:assurance}) gives the complete answer.
\end{problem}

\section{Compression as the through-line}

Gathering the whole paper together yields a single thread: \term{compression}.

\begin{itemize}
\item Complex systems compress themselves by \emph{levels}, so that every large system
      grows a kernel internal to its domain (Chapter \ref{ch:infrastructure});
\item Bounded intelligence compresses information by \emph{comprehension}, so that a
      finite cognition accommodates an unbounded world (Chapter \ref{ch:comprehension});
\item Engineering compresses the cost of change by \emph{architecture}, so that a local
      modification does not require global comprehension (Chapter \ref{ch:design});
\item Civilisation compresses trust by \emph{rules}, so that a million lines of code of
      unknown provenance can be run safely (Chapter \ref{ch:intent}).
\end{itemize}

The mechanism of compression is the same in all four cases: quotienting by the kernel of
the abstraction map (Proposition \ref{prop:oblivious}). \emph{All the information a level
exposes to the level above is the quotient of lower-level behaviour by that kernel.}

\section{The core proposition}

If the whole paper is to be compressed into one testable sentence, this paper chooses the
following:

\begin{principle}[Core proposition]\label{pr:core}
A complex system is comprehensible to a bounded intelligence if and only if it is
compressible; and it is compressible if and only if it admits a separation of scales in
which every level has an explicit, declared \emph{oblivious set}.
\end{principle}

Three consequences follow directly.

First, it turns ``comprehensibility'' from a subjective impression into a structural
property: comprehensibility is equivalent to the existence of a quotient map under which
task-relevant predictive power is preserved. This makes comprehensibility an
\emph{objective function} of design rather than a retrospective assessment.

Second, it redefines ``architecture'' as \emph{manufactured separation of scales}. Nature
obtains levels free from separation of scales and renormalisation; software levels do not
emerge spontaneously and must be paid for in interfaces, encapsulation and invariants,
and they are liable to leak. The core deliverable of architecture is therefore not a
module list but the \emph{declaration of the oblivious set of every level}.

Third, it gives a concrete rule of action for an AI: do not attempt to exhaust low-level
information; maintain a layered, falsifiable compression model that counterexamples can
update; obtain gradient by probing rather than by reading; and take task sufficiency as
the stopping condition.

\section{Closing remarks}

The paper begins with three engineering intuitions and concludes with a joint proposition
about structure and assurance. What is established is not the answers but \emph{which
questions have no answer}, and why:

\begin{itemize}
\item the gap between intent and specification has no mechanical solution, and its
      grounds are epistemological rather than logical; it is in substance the frame
      problem, and is therefore avoidable only in interaction;
\item the basis of trust for mechanical assurance can be narrowed but not eliminated;
      and ``verification is cheap'' must be decomposed into two independent propositions,
      cheap checking time and cheap certificate size, of which the latter is the usual
      bottleneck;
\item the reliability of consensus assurance is dominated by correlated error modes, so
      that \emph{diversity} rather than \emph{number} is the design variable, and
      diversity is ultimately a problem of incentive design, not of statistics;
\item the structural convergence of large systems is probabilistic, domain-specific and
      shaped by organisational constraints, rather than unique;
\item the duty of architecture is not to lower the cost of a single change but to lower
      the \emph{exponent} with which that cost grows with size.
\end{itemize}

Set side by side, these items share one shape: \emph{complexity compression is a
necessity driven by size, and the form the compression takes is a contingency driven by
domain; while the boundary of assurance falls exactly between what a machine can decide
and what must be validated by the party that bears the consequences.}

The price is the surrender of a picture expressible as a single theorem; the gain is that
every claim sits where it can be tested. For a framework about how to establish
trustworthy conclusions on top of untrusted complexity, this is the only self-consistent
ending: \emph{it must pass its own test.}
